
\begin{TP}[La formule de Pick]

On va s'intéresser au théorème de Pick qui permet de calculer l'aire d'un polygone construit sur du papier pointé et dont les sommets sont des points du papier. Voici, par exemple, un tel polygone :
\begin{center} \includegraphics[width=3.8cm]{Pick} \end{center}

\partie{Georg Alexander Pick}
Par groupe, en vous documentant, répondez aux questions suivantes.
\begin{enumerate}
 \item Où et à quelle époque, Georg Alexander Pick a‑t‑il vécu ?
 \item Quels sont les domaines dans lesquels Georg Alexander Pick a travaillé ?
 \item Faites la synthèse des réponses de chaque groupe.
\end{enumerate}

\partie{Formule de Pick}
Pour un polygone construit sur du papier pointé et dont les sommets sont des points du papier, on appelle $N$ le nombre de points situés sur son contour et $P$ le nombre de points situés à l'intérieur. Le théorème de Pick donne la formule pour calculer l'aire $\mathcal{A}$ de ce polygone :
\begin{center} $\mathcal{A} = 0,5 \cdot N + P - 1$ ; l'unité est le carreau. \end{center}
%\vspace{0.5cm}
\begin{enumerate}
 \item Chaque groupe calcule l'aire du rectangle ci‑dessous en utilisant la formule habituelle puis en utilisant la formule de Pick. Comparez avec les réponses des autres groupes. \\[1em]
\begin{center} \includegraphics[width=2.9cm]{Pick2} \end{center}

 \item Chaque groupe construit cinq polygones sur du papier pointé, avec chaque sommet placé sur un point.
 \item Échangez ensuite avec un autre groupe les polygones. Calculez l'aire de chacun des polygones reçus.
 \end{enumerate}
\end{TP}

